\chapter{Methodology}
\label{chap:methodology}

%====================================================
\section{Overview of the Methodology}

This chapter presents a comprehensive methodological framework developed for explainable machine learning–based fault classification of power transformers using Dissolved Gas Analysis (DGA). The methodology is structured in a systematic and sequential manner to ensure that each stage of the process contributes directly to the overall research objectives. The primary aim of this methodology is to ensure reproducibility, consistency, and interpretability throughout the entire machine learning pipeline.

The methodological design is carefully structured to support three key research objectives, which include the evaluation of feature sets, the comparison of traditional machine learning algorithms, and the interpretation of model decisions using explainable artificial intelligence techniques. Each of these objectives is addressed through a dedicated stage within the pipeline, ensuring that the entire workflow remains aligned and logically structured \parencite{Dladla2025SLR}.

The importance of a structured methodology lies in ensuring that each step in the data processing and modelling pipeline contributes meaningfully to the final outcome. In machine learning-based diagnostic systems, especially in engineering applications such as transformer fault classification, a well-defined methodological framework is essential for ensuring reliability and consistency across experiments \parencite{Cao2024Advancement}.

Furthermore, the pipeline is designed in a modular format, where each stage operates independently but contributes cumulatively to the final classification and interpretation process. This modularity ensures that any modification in one stage does not compromise the integrity of the entire system, thereby improving robustness and scalability of the proposed approach.

The chapter is organized in a logical progression that begins with raw data processing and continues through feature extraction, feature selection, model training, and evaluation, before concluding with explainable artificial intelligence interpretation. This sequential structure ensures a clear and traceable flow of information throughout the methodology.

%====================================================
\section{Data Ingestion and Initial Cleaning}

Data ingestion represents the initial stage of the machine learning pipeline, where benchmark DGA datasets are collected and loaded into the computational environment for analysis. These datasets consist of transformer oil dissolved gas concentrations and their corresponding fault classification labels. The integrity of this stage is critical because any inconsistencies in raw data may directly influence the performance of subsequent machine learning models.

The dataset includes key dissolved gases such as hydrogen (H$_2$), methane (CH$_4$), ethane (C$_2$H$_6$), ethylene (C$_2$H$_4$), acetylene (C$_2$H$_2$), carbon monoxide (CO), and carbon dioxide (CO$_2$). These gases are widely used in transformer diagnostics because they are directly associated with different fault mechanisms occurring within transformer insulation systems \parencite{Ali2023ConventionalDGA}.

During data ingestion, the dataset structure is carefully verified to ensure that all required attributes are present and correctly formatted. This includes checking the consistency of column names, data types, and label distributions. Any structural inconsistencies are corrected at this stage to prevent downstream errors during model training.

Initial cleaning is then performed to address data quality issues. Missing values are identified and treated using appropriate statistical or imputation methods depending on the severity of missingness. This ensures that the dataset remains complete without introducing bias into the learning process.

In addition, noise filtering is applied to remove inconsistent or outlier values that may distort the learning patterns of machine learning models. This step is important because real-world industrial datasets often contain measurement errors or abnormal readings that do not represent actual transformer operating conditions.

Normalization is also applied as part of preprocessing to ensure that all gas concentration values are represented on a comparable scale. This is particularly important because different gases exhibit different magnitude ranges, and unscaled data may lead to biased learning in distance-sensitive algorithms.

%====================================================
\section{Feature Extraction Framework}

Feature extraction is a fundamental stage in the machine learning pipeline where raw DGA measurements are transformed into meaningful representations that can be effectively processed by classification models. The purpose of this stage is to convert raw numerical gas values into structured feature representations that capture relevant patterns associated with transformer faults.

In transformer diagnostics, raw gas concentrations alone may not fully represent the underlying relationships between fault conditions. Therefore, feature extraction plays a critical role in enhancing the representational power of the dataset by introducing different perspectives of the same underlying data structure \parencite{Hu2025InterpretableEdge}.

This framework is designed to extract multiple types of features, each focusing on a different aspect of the data. These include density-based representations, transformation-based representations, structural representations, and statistical representations. Each feature type contributes uniquely to the overall understanding of transformer fault behavior.

%====================================================
\subsection{Density Features}

Density features represent the distribution characteristics of gas concentrations within the dataset. This feature type focuses on how frequently specific gas concentration ranges occur and how data points are distributed across different fault categories.

The purpose of density-based representation is to capture the overall distribution pattern of gases rather than focusing solely on individual values. This allows the model to understand how gas concentrations are distributed across different operating conditions of transformers.

By analyzing density patterns, it becomes possible to identify regions of high and low concentration activity, which may correspond to different types of fault conditions. This representation provides an additional layer of information beyond raw numerical values.

%====================================================
\subsection{Radon Transform Features}

Radon transform features are derived through mathematical transformations that project data into alternative representation spaces. The purpose of this transformation is to capture hidden structural information that may not be directly observable in the original feature space.

This transformation technique is useful for identifying global patterns within the dataset by analyzing how data behaves under different projection angles. It allows the model to capture structural relationships that may otherwise remain hidden in conventional feature representations.

The use of transformation-based features enhances the ability of the machine learning model to learn more complex relationships within the dataset without altering the original data semantics.

%====================================================
\subsection{Geometric Features}

Geometric features describe the structural relationships between gas variables in a spatial context. These features are designed to capture the shape and configuration of data points when represented in multidimensional space.

The main objective of geometric representation is to understand how different gas variables interact with each other in terms of spatial distribution. This includes analyzing how gas combinations form specific patterns that may correspond to different fault types.

By representing data geometrically, the model gains a better understanding of relational structures within the dataset, which can improve classification performance in complex scenarios.

%====================================================
\subsection{Statistical Features}

Statistical features represent the numerical properties of gas distributions using standard statistical measures. These include mean, variance, skewness, and kurtosis, which collectively describe the central tendency, dispersion, and shape of the data distribution.

The purpose of statistical representation is to summarize the behavior of gas concentrations in a compact mathematical form. This allows the model to capture essential characteristics of the dataset without relying on raw values alone.

These features are particularly useful in identifying variability patterns and distribution asymmetry, which may be indicative of specific transformer fault conditions.

%====================================================
\section{Feature Expansion and Selection}

Feature expansion and selection represent a critical stage in the methodology where the feature space is enhanced and refined to improve model performance. Feature expansion increases the diversity of feature representations, while feature selection ensures that only the most relevant features are retained.

The importance of this stage lies in balancing model complexity and predictive performance. Excessive features may introduce noise and redundancy, while insufficient features may lead to loss of important information \parencite{Laburu2024Insights}.

%====================================================
\subsection{Feature Expansion}

Feature expansion refers to the process of generating additional feature representations from existing gas measurements. This is achieved through mathematical transformations and combinations of original features.

The purpose of feature expansion is to enrich the dataset by introducing new perspectives of the same underlying information. This allows machine learning models to capture more complex relationships between gas variables and fault conditions.

By expanding the feature space, the model gains access to a more expressive representation of the data, which may improve classification capability.

%====================================================
\subsection{Feature Selection Strategy}

Feature selection is the process of identifying and retaining only the most relevant features for model training. This step is essential for reducing redundancy and improving computational efficiency.

The selection process ensures that irrelevant or weakly correlated features are removed, allowing the model to focus on the most informative variables. This contributes to improved generalization performance and reduces the risk of overfitting.

%====================================================
\subsection{Stage 1: Global Pre-Filtering}

Global pre-filtering is the initial stage of feature selection where features with minimal contribution are removed. This step is performed to reduce noise and simplify the feature space before more detailed selection is applied.

By eliminating weak features early, the dataset becomes more structured and easier to analyze in subsequent stages.

%====================================================
\subsection{Stage 2: Fine Selection Tracks}

Fine selection tracks involve a more detailed evaluation of the remaining features. This stage focuses on identifying the most optimal subset of features that contribute significantly to classification performance.

The objective of this stage is to refine the feature set to ensure that only the most relevant and informative features are retained for model training.

%====================================================
\section{Data Splitting and Hybrid Balancing}

Data splitting is performed to divide the dataset into training and testing subsets. This ensures that model performance can be evaluated on unseen data, providing a reliable measure of generalization capability.

Cross-validation is also applied to further enhance the robustness of the evaluation process. This helps reduce variability in performance estimates and ensures that the model is not overfitted to a specific data split.

Class imbalance is addressed using hybrid balancing techniques such as SMOTE. This ensures that minority fault classes are adequately represented during model training, preventing bias toward majority classes \parencite{Azmi2025Imbalance}.

%====================================================
\section{Model Training and Evaluation}

Model training involves the application of machine learning algorithms to the processed feature sets for fault classification. The training process is conducted in a controlled environment to ensure fairness and consistency across all models.

Evaluation is performed using standard classification metrics, including accuracy, precision, recall, and F1-score. These metrics provide a comprehensive assessment of model performance, especially in scenarios involving class imbalance.

The results of the evaluation are used to compare the effectiveness of different feature sets and machine learning models in DGA-based transformer fault classification.